% Part XXVII: The W(3,3) graph Riemann hypothesis as a theorem
% \input'd into w33_paper.tex after Part XXVI

\part*{Part XXVII\,---\,The $W(3,3)$ Graph Riemann Hypothesis as a Theorem}
\addcontentsline{toc}{part}{Part XXVII\,---\,The $W(3,3)$ Graph Riemann Hypothesis as a Theorem}

\bigskip
\noindent\textbf{Overview.}\;
Supplement~H verified numerically that all non-trivial poles of the
Ihara zeta function $\zeta_{W(3,3)}(u)$ lie on the circle
$|u|=1/\sqrt{k-1}=1/\sqrt{11}$.
Part~XXVII converts that verification into a proof from the exact
adjacency spectrum $12^1,2^{24},(-4)^{15}$.  The result is a complete
finite graph-RH theorem.  It is an analogue of the classical Riemann
Hypothesis, not a proof of the classical statement; an explicit transfer
theorem to the completed Riemann zeta function remains necessary.

%-----------------------------------------------------------------------
\section{The Ihara Zeta Function of $W(3,3)$}
\label{sec:ihara}

\begin{definition}[Ihara zeta function]
Let $G$ be a finite $k$-regular graph with adjacency eigenvalues
$k=\mu_0\ge\mu_1\ge\cdots\ge\mu_{v-1}$.  Its Ihara zeta function obeys
\begin{equation}
  \zeta_G(u)^{-1}
  = (1-u^2)^{E-v}\prod_{j=0}^{v-1}
    \bigl(1-\mu_j u+(k-1)u^2\bigr).
  \label{eq:ihara}
\end{equation}
\end{definition}

\begin{proposition}[Ihara determinant for $W(3,3)$]
\label{prop:ihara_det}
For $W(3,3)$ with $v=40$, $E=240$, $k=12$, and spectrum
$12^1,2^{24},(-4)^{15}$,
\begin{equation}
  \zeta_{W(3,3)}(u)^{-1}
  = (1-u^2)^{200}
    (1-u)(1-11u)
    (1-2u+11u^2)^{24}
    (1+4u+11u^2)^{15}.
  \label{eq:ihara_w33}
\end{equation}
The non-trivial poles are the roots of the last two quadratic factors.
\end{proposition}

\begin{proof}
Since $E-v=240-40=200$, Bass's determinant formula gives the stated
prefactor.  The eigenvalue $12$ contributes
$1-12u+11u^2=(1-u)(1-11u)$, while the restricted eigenvalues $2$ and
$-4$ contribute the two non-trivial factors with multiplicities $24$
and $15$.
\end{proof}

%-----------------------------------------------------------------------
\section{The Ramanujan Bound}
\label{sec:ramanujan}

\begin{definition}[Ramanujan graph]
A $k$-regular graph is Ramanujan if every non-trivial adjacency
eigenvalue $\mu$ satisfies $|\mu|\le 2\sqrt{k-1}$.
\end{definition}

\begin{theorem}[Ramanujan bound for $W(3,3)$]
\label{thm:ramanujan}
The graph $W(3,3)$ is Ramanujan.  Its non-trivial eigenvalues satisfy
\begin{equation}
  |2|<2\sqrt{11},
  \qquad
  |-4|<2\sqrt{11}.
  \label{eq:ram_bound}
\end{equation}
\end{theorem}

\begin{proof}
The exact restricted spectrum is $\{2,-4\}$ and
$2\sqrt{k-1}=2\sqrt{11}\approx6.633$.  Direct comparison proves both
inequalities.
\end{proof}

%-----------------------------------------------------------------------
\section{Compatible $H_4$ Spectral Scale}
\label{sec:h4_comparison}

\begin{proposition}[$H_4$ comparison, not implication]
The four $H_4$ Coxeter values used in Part~XXIV,
\begin{equation}
  \varphi+1,\quad \varphi,\quad \varphi-1,\quad 2-\varphi,
\end{equation}
all lie strictly inside $(-2\sqrt{11},2\sqrt{11})$.
\end{proposition}

\begin{proof}
Their approximate values are $2.618,1.618,0.618,0.382$, all smaller in
absolute value than $2\sqrt{11}\approx6.633$.
\end{proof}

\begin{remark}
This is a compatible scale comparison only.  It does not force the
Ramanujan bound, and arbitrary $\mathbb{Q}(\varphi)$-linear combinations
are not bounded by the Ramanujan interval.  The graph theorem rests on
the exact $W(3,3)$ adjacency spectrum.
\end{remark}

%-----------------------------------------------------------------------
\section{The Graph Riemann Hypothesis}
\label{sec:grh_theorem}

\begin{theorem}[Graph Riemann Hypothesis for $W(3,3)$]
\label{thm:grh}
All non-trivial poles of $\zeta_{W(3,3)}(u)$ lie on
\begin{equation}
  |u|=\frac{1}{\sqrt{k-1}}=\frac{1}{\sqrt{11}}.
  \label{eq:grh_circle}
\end{equation}
\end{theorem}

\begin{proof}
For $\mu\in\{2,-4\}$, a non-trivial pole is a root of
$1-\mu u+11u^2=0$.  The discriminant $\mu^2-44$ is negative, so the two
roots are complex conjugates.  Their product is $1/11$, hence each root
has modulus $1/\sqrt{11}$.
\end{proof}

\begin{remark}[Role of the selector $q!=2q$]
The identity $q!=2q$ uniquely selects $q=3$ within the repository's
architectural language.  It does not logically imply the Ramanujan bound
or the graph-RH.  Those statements follow from the exact spectrum of the
selected graph.
\end{remark}

%-----------------------------------------------------------------------
\section{Reciprocity and Critical-Line Normalization}
\label{sec:functional}

\begin{proposition}[Quadratic reciprocity]
For $q_0=k-1=11$ and
$P_\mu(u)=1-\mu u+q_0u^2$,
\begin{equation}
  P_\mu\!\left(\frac{1}{q_0u}\right)
  =\frac{P_\mu(u)}{q_0u^2}.
\end{equation}
Thus roots occur in reciprocal pairs $u\leftrightarrow1/(11u)$.
\end{proposition}

\begin{proof}
Substitution and multiplication by $11u^2$ give
$11u^2-\mu u+1=P_\mu(u)$.
\end{proof}

\begin{proposition}[Correct exponential coordinate]
Under
\begin{equation}
  u=11^{-s},
\end{equation}
the graph critical circle $|u|=11^{-1/2}$ is equivalent to
$\operatorname{Re}(s)=1/2$.
\end{proposition}

\begin{proof}
Taking absolute values gives
$|u|=11^{-\operatorname{Re}(s)}$.  Equating this with $11^{-1/2}$ yields
$\operatorname{Re}(s)=1/2$.
\end{proof}

\begin{remark}[Normalization correction]
If one instead writes $u=(\sqrt{11})^{-s}$, then
$|u|=1/\sqrt{11}$ corresponds to $\operatorname{Re}(s)=1$, not $1/2$.
The base-$11$ coordinate is the correct normalization for the stated
critical-line analogy.
\end{remark}

%-----------------------------------------------------------------------
\section{Connection to the Classical Riemann Hypothesis}
\label{sec:crh}

\begin{remark}
The classical Riemann Hypothesis concerns the zeros of the completed
Riemann zeta function.  Theorem~\ref{thm:grh} is a finite-graph analogue:
Ramanujan adjacency bounds are equivalent to Ihara poles lying on the
critical circle.  No equivalence between $\zeta_{W(3,3)}(u)$ and the
classical completed zeta function is proved here.  A Clay-level argument
would require a rigorous determinant, trace-formula, automorphic, or
limiting transfer that identifies the classical zero set.
\end{remark}

%-----------------------------------------------------------------------
\section{Computable Residue Frontier}
\label{sec:p18}

\begin{enumerate}[label=\textbf{P\arabic*.},start=18]
  \item \textbf{Graph-zeta pole residues.}
        Compute the residues at all $24+15=39$ non-trivial pole slots and
        determine their exact number fields and denominator ideals.  This
        is a finite computer-algebra problem and should be reported as a
        certificate, not assumed from multiplicity counts alone.
\end{enumerate}
